\documentclass[11pt]{article}

\usepackage[a4paper,margin=1in]{geometry}
\usepackage[T1]{fontenc}
\usepackage[utf8]{inputenc}
\usepackage{lmodern}
\usepackage{amsmath,amssymb,mathtools}
\usepackage{slashed}
\usepackage{booktabs}
\usepackage{float}
\usepackage{array}
\usepackage{graphicx}
\usepackage{xcolor}
\usepackage{hyperref}

\hypersetup{
  colorlinks=true,
  linkcolor=blue!60!black,
  citecolor=blue!60!black,
  urlcolor=blue!60!black
}

\newcommand{\ii}{\mathrm{i}}
\newcommand{\Ds}{D_s}
\newcommand{\Dsone}{D_{s1}}
\newcommand{\Bs}{B_s}
\newcommand{\Bsone}{B_{s1}}
\newcommand{\GeV}{\mathrm{GeV}}
\newcommand{\MeV}{\mathrm{MeV}}
\newcommand{\keV}{\mathrm{keV}}
\newcommand{\qqbar}{\langle\bar q q\rangle}
\newcommand{\ssbar}{\langle\bar s s\rangle}
\newcolumntype{L}[1]{>{\raggedright\arraybackslash}p{#1}}

\title{Radiative Transitions of Strange Heavy Axial Mesons from Light-Cone QCD Sum Rules}
\author{S. Bilmis et al.}
\date{\today}

\begin{document}
\maketitle

\begin{abstract}
We present a light-cone QCD sum-rule analysis of
$\Dsone(2460)\to\Ds\gamma$, $\Dsone(2536)\to\Ds\gamma$, and
$\Bsone(5830)\to\Bs\gamma$.  The calculation keeps explicit strange-quark-mass
terms and includes hard photon emission from both quark lines, two-particle
photon distribution amplitudes, and the three-particle photon-amplitude
corrections induced by the background-gluon term in the quark propagator.  We
separate the basis-current sum rules from the physical-current residue
normalization.  For the charm-strange channel the physical residues are
projected from an $AA/AB/BB$ two-point matrix at the previously determined
mixing angle; the bottom-strange residues are obtained from the direct
heavy-flavour continuation of the same matrix calculation.  Using the modern
lattice normalization of the
leading-twist photon distribution amplitude, we find
$\Gamma[\Dsone(2460)\to\Ds\gamma]=34.96[28.06,43.53]~\keV$,
$\Gamma[\Dsone(2536)\to\Ds\gamma]=0.988[0.092,3.555]~\keV$, and
$\Gamma[\Bsone(5830)\to\Bs\gamma]=8.702[6.474,11.23]~\keV$.
The operator-weighted gluonic and electromagnetic convolutions are evaluated
with the exact off-shell double-Borel numerator reduction and agree
term-by-term between the Python and Mathematica implementations.
\end{abstract}

\section{Introduction}

Radiative transitions of positive-parity strange heavy mesons probe the
interplay of heavy-quark spin symmetry, light-quark dynamics, and threshold
effects.  The two charm-strange axial states, $\Dsone(2460)$ and
$\Dsone(2536)$, are especially useful because the lower state has a measured
radiative branching fraction \cite{BaBar2006}, while the upper state has a
measured total width but no prominent $\Ds\gamma$ signal \cite{BaBar2011}.  The
bottom-strange state $\Bsone(5830)$ is observed in hadronic spectroscopy
\cite{CMS2018}; its radiative decay has not yet been measured.

The aim of this paper is not only to quote updated numbers, but also to make the
normalization of the calculation unambiguous.  The light-cone OPE gives
transition form factors for basis currents.  The conversion to the physical
$1^+$ states requires the residues of the physical mixed currents.  We therefore
organize the analysis in stages: first compute the axial and tensor basis
amplitudes, then add the photon-DA corrections, and finally apply the
physical-current normalization.

The hierarchy between $\Dsone(2460)\to\Ds\gamma$ and
$\Dsone(2536)\to\Ds\gamma$ was recently emphasized by Bondar and by
Bondar and Milstein \cite{Bondar2025Why,BondarMilstein2025}.  Their
quark-model explanation is that the physical $J=1$ states are mixed spin
configurations and that the E1 amplitudes enter with different relative signs.
Below we show that the same mechanism appears directly in the LCSR
decomposition of the physical amplitudes.

Two photon-input scenarios are considered throughout.  The first is the legacy
magnetic-susceptibility normalization of the photon distribution amplitudes
\cite{Ball:2002ps,Rohrwild2007}.  The second uses the recent lattice
normalization of the leading-twist photon distribution amplitude
\cite{Bacchio2024}.  The latter is our preferred modern scenario; the former is
kept to show continuity with older light-cone sum-rule analyses.

\section{Hadronic Conventions}

We use $p'=p+q$, where $p'$ is the momentum of the initial axial meson, $p$ is the
momentum of the final pseudoscalar meson, and $q$ is the real-photon momentum.
The pseudoscalar interpolating current is
\begin{equation}
  J_5=\bar s\,\ii\gamma_5 Q,\qquad Q=c,b .
\end{equation}
For the axial channel we introduce two basis currents,
\begin{align}
  J^A_\mu &= \bar s\gamma_\mu\gamma_5 Q,\\
  J^B_\mu &= \frac{\ii}{m_Q+m_s}\bar s\sigma_{\mu\alpha}p'^\alpha\gamma_5 Q.
  \label{eq:basis-currents}
\end{align}
The momentum in the tensor current is the initial-state momentum $p'^\alpha$.
The physical-current combinations are
\begin{align}
  J_{1\mu} &= \sin\theta\,J^A_\mu+\cos\theta\,J^B_\mu,\\
  J_{2\mu} &= \cos\theta\,J^A_\mu-\sin\theta\,J^B_\mu ,
  \label{eq:physical-currents}
\end{align}
We use the sector-specific results of our dedicated mixing-angle study
\cite{AlievBilmisSavci2024},
\begin{equation}
\theta_{D_s}=26.6\pm0.6^\circ,\qquad
\theta_{B_s}=38.5\pm0.1^\circ.
\end{equation}
The heavy-quark-limit value $35.3^\circ$ and the interval
$25^\circ\leq\theta\leq45^\circ$ are retained only as benchmark and
sensitivity tests; the latter is not interpreted as a statistical confidence
interval.

The radiative matrix element is parametrized as
\begin{equation}
  \langle P(p)\gamma(q,\varepsilon)|A(p',\eta)\rangle
  =
  eG\left[(\eta\cdot\varepsilon^*)(p\cdot q)
  -(\eta\cdot q)(p\cdot\varepsilon^*)\right].
  \label{eq:e1-definition}
\end{equation}
The corresponding width is
\begin{equation}
  \Gamma(A\to P\gamma)
  =
  \frac{\alpha_{\rm em}}{3}\,G^2
  \left(\frac{m_A^2-m_P^2}{2m_A}\right)^3 .
  \label{eq:width}
\end{equation}

\section{Light-Cone Sum Rule}

The starting correlation function is
\begin{equation}
  \Pi_\mu(p,q)=\ii\int d^4x\,e^{\ii p\cdot x}
  \langle\gamma(q,\varepsilon)|
  T\{\bar s(x)\Gamma_\mu Q(x)\,\bar Q(0)\ii\gamma_5 s(0)\}
  |0\rangle ,
  \label{eq:correlator}
\end{equation}
where $\Gamma_\mu$ is either $J^A_\mu$ or $J^B_\mu$.  We project onto the
gauge-invariant E1 tensor
\begin{equation}
  T_{\mu\nu}=(p\cdot q)g_{\mu\nu}-q_\mu p_\nu .
\end{equation}

The OPE is constructed with quark propagators in an external electromagnetic
field.  Both hard photon emission and soft photon distribution-amplitude terms
are retained.  In addition, the background-gluon insertion in the propagator
generates the three-particle photon-DA terms.  The ordinary vacuum-condensate
pieces of the propagator are not inserted as separate three-point corrections:
in the photon LCSR they enter through photon matrix elements and normalizations
such as $\chi_s\ssbar$.  This is separate from the local two-point sum rules
used to determine decay constants, where ordinary SVZ condensates are part of
the OPE.

After Borel transformation and continuum subtraction the basis-current
amplitudes are denoted $G_A$ and $G_B$.  The charm physical couplings are then
formed as
\begin{align}
  G_{2460} &=
  \frac{\sin\theta\,f_A\,G_A+\cos\theta\,f_B\,G_B}{f_1},\\
  G_{2536} &=
  \frac{\cos\theta\,f_A\,G_A-\sin\theta\,f_B\,G_B}{f_2}.
  \label{eq:stage3}
\end{align}
Here $f_1$ and $f_2$ are obtained by projecting the normalized-current
$AA/AB/BB$ two-point QCD sum-rule matrix at the external $D_s$ angle.  The
LO+$d=3$+$d=5$ Monte Carlo gives
\begin{equation}
f_1=0.4046[0.3799,0.4291]~\GeV,\qquad
f_2=0.1677[0.1594,0.1756]~\GeV.
\end{equation}
Diagonalizing this truncated matrix gives
$30.63[29.94,31.29]^\circ$ only as a diagnostic.  The normalized off-diagonal
residual after the external-angle projection is $0.183[0.142,0.225]$.

Wang's current arXiv v4 result is
$f_A(D_{s1})=0.345\pm0.017~\GeV$ \cite{Wang2015}; the former
$0.245\pm0.017~\GeV$ entry was explicitly corrected.  Reproducing that number
is neither necessary nor sufficient to verify the mixed residues: Wang uses a
pure $J_A$ current and a one-pole ansatz, whereas our $AA$ channel contains the
two physical poles after the mixed-current projection.  We therefore make the
comparison at two levels.  The direct LO+$d=3$+$d=5$ $AA$ OPE evaluated in
Wang's window gives $0.2686[0.2607,0.2768]~\GeV$.  Independently, projecting
our two physical poles back onto $J_A$ gives a Borel-weighted one-pole-equivalent
residue $0.2303[0.2194,0.2418]~\GeV$.  The analytic projection is written in
Appendix~\ref{app:two-point}.  Both numbers are normalization diagnostics;
neither is used as an input or presented as a reproduction of Wang's complete
OPE.

For the bottom-strange channel we evaluate the same normalized-current
$AA/AB/BB$ OPE with $m_Q=m_b$ and project it at the previous-study angle
$\theta_{B_s}=38.5\pm0.1^\circ$.  In the selected two-point window
$M^2=6.50$--$7.30~\GeV^2$, separate thresholds are fitted to the model lower
pole and the observed $\Bsone(5830)$.  The direct Monte Carlo gives
\begin{equation}
 f_1^{(B_s)}=0.5358[0.5096,0.5586]~\GeV,\qquad
 f_2^{(B_s)}=0.08896[0.08496,0.09278]~\GeV .
\end{equation}
Of 2000 generated samples, 1757 satisfy all stated stability criteria.
No Pullin--Zwicky decay constant and no
off-diagonal overlap closure enters the final bottom normalization; their
earlier use is retained only in the internal historical comparison.

\section{Numerical Inputs and Scale Convention}

Table~\ref{tab:inputs-polished} lists the main inputs.  Hadron masses, quark
masses, charges, and $\alpha_{\rm em}$ are taken from the PDG review and
online listings \cite{PDG2024,PDG2026}.  The pseudoscalar decay constants are
taken in the lattice context summarized by the published FLAG review
\cite{FLAG2024}.  Photon DAs follow Ball, Braun and
Kivel \cite{Ball:2002ps}; the legacy susceptibility comparison follows
Rohrwild \cite{Rohrwild2007}; the modern photon normalization follows the
lattice result of Bacchio et al. \cite{Bacchio2024}.  The input bookkeeping is
kept in \texttt{inputs\_table.csv} and in the generated citation map
\texttt{outputs/input\_citation\_map.csv}.

The photon-DA normalizations are quoted at $\mu=1~\GeV$, while the quark masses
are the usual $\overline{\rm MS}$ inputs, e.g. $m_s(2~\GeV)$ and $m_c(m_c)$.
This is a standard phenomenological convention rather than a fully matched
single-scale setup.  We therefore treat the lattice photon-normalized scenario
as the preferred modern input and retain the susceptibility scenario as a
comparison with older analyses.

\begin{table}[htbp]
  \centering
  \resizebox{\textwidth}{!}{%
  \begin{tabular}{llll}
    \toprule
    Quantity & Symbol & Value & Source/use\\
    \midrule
    quark masses & $m_c,m_b,m_s$ & $1.27(2),\,4.18(3),\,0.093(11)~\GeV$
      & PDG; explicit $m_s$ terms kept\\
    charm masses & $M_{2460},M_{2536},m_{\Ds}$
      & $2.4595,\,2.53511,\,1.96835~\GeV$ & PDG\\
    bottom masses & $M_{5830},m_{\Bs}$
      & $5.82870,\,5.36692~\GeV$ & PDG/CMS\\
    lower $\Bsone$ diagnostic mass & $M_{\Bsone}^{\rm low}$
      & $5.750(26)~\GeV$ & lattice-inspired diagnostic \cite{Lang2015}\\
    pseudoscalar decay constants & $f_{\Ds},f_{\Bs}$
      & $0.2499,\,0.2303~\GeV$ & FLAG/lattice context\\
    charm mixed residues & $f_1,f_2$
      & $0.4046[0.3799,0.4291],\,0.1677[0.1594,0.1756]~\GeV$
      & fixed-angle $AA/AB/BB$ projection\\
    pure axial benchmark & $f_A(D_{s1})$
      & $0.345(17)~\GeV$ & corrected Wang v4 pure-current check\\
    bottom physical constants & $f_1^{(B_s)},f_2^{(B_s)}$
      & $0.5358[0.5096,0.5586],\,0.08896[0.08496,0.09278]~\GeV$
      & direct $AA/AB/BB$ projection\\
    bottom two-point window & $M_{\rm 2pt}^2$
      & $6.50$--$7.30~\GeV^2$ & pole/OPE/mass stability\\
    legacy photon input & $\chi_s(1~\GeV)$
      & $+3.15(30)~\GeV^{-2}$ & comparison scenario\\
    lattice photon input & $f_{\gamma,s}^{\perp}(1~\GeV)$
      & $-55.1(1.9)~\MeV$ & preferred scenario\\
    \bottomrule
  \end{tabular}%
  }
  \caption{Main numerical inputs used in the analysis.}
  \label{tab:inputs-polished}
\end{table}

\begin{table}[htbp]
  \centering
  \resizebox{\textwidth}{!}{%
  \begin{tabular}{lllll}
    \toprule
    Quantity & Current & Central value & Interval & Interpretation\\
    \midrule
    Wang v4 $f_A$ & pure $J_A$, one pole & $0.345$ & $\pm0.017$ & external benchmark\\
    our direct $AA$ OPE & pure $J_A$ & $0.2686$ & $[0.2607,0.2768]$ & truncation benchmark\\
    projected $AA$ equivalent & $J_A$, two poles & $0.2303$ & $[0.2194,0.2418]$ & one-pole-equivalent diagnostic\\
    $f_1$ & mixed $J_1$ & $0.4046$ & $[0.3799,0.4291]$ & not directly Wang-comparable\\
    $f_2$ & mixed $J_2$ & $0.1677$ & $[0.1594,0.1756]$ & not directly Wang-comparable\\
    $\sin\theta f_1$ & $J_A\to D_{s1}(2460)$ & $0.1812$ & $[0.1692,0.1930]$ & projection diagnostic\\
    $\cos\theta f_2$ & $J_A\to D_{s1}(2536)$ & $0.1498$ & $[0.1424,0.1573]$ & projection diagnostic\\
    \bottomrule
  \end{tabular}%
  }
  \caption{Pure-current and projected-current comparison.  Wang's number,
  our direct truncated $AA$ OPE, and the projected two-pole equivalent do not
  share the same OPE truncation and pole model.  The table is therefore a
  normalization benchmark, not a direct validation of either mixed residue.
  Decay constants are in GeV.}
  \label{tab:residue-benchmark-polished}
\end{table}

\section{Results}

The final numerical results are summarized in
Table~\ref{tab:results-polished}.  We
quote medians and 16--84 percentile intervals rather than Gaussian means and
standard deviations.  This is important because the $\Dsone(2536)$ and
$\Bsone(5830)$ amplitudes can pass close to zero; their width distributions are
therefore non-Gaussian.

\begin{table}[htbp]
  \centering
  \resizebox{\textwidth}{!}{%
  \begin{tabular}{lllllcc}
    \toprule
    Sector & Decay & Photon input & Mixing & Normalization
      & $G~(\GeV^{-1})$ & $\Gamma~(\keV)$\\
    \midrule
    $D_s$ & $\Dsone(2460)\to\Ds\gamma$ & legacy $\chi_s$
      & $26.6(6)^\circ$ & $f_1=0.405[0.381,0.429]$
      & $-0.542[-0.611,-0.479]$ & $61.66[48.37,78.46]$\\
    $D_s$ & $\Dsone(2536)\to\Ds\gamma$ & legacy $\chi_s$
      & $26.6(6)^\circ$ & $f_2=0.168[0.159,0.175]$
      & $-0.0790[-0.153,-0.00266]$ & $2.220[0.226,7.379]$\\
    $D_s$ & $\Dsone(2460)\to\Ds\gamma$ & lattice $f_{\gamma,s}^{\perp}$
      & $26.6(6)^\circ$ & $f_1=0.405[0.379,0.430]$
      & $-0.408[-0.455,-0.365]$ & $34.96[28.06,43.53]$\\
    $D_s$ & $\Dsone(2536)\to\Ds\gamma$ & lattice $f_{\gamma,s}^{\perp}$
      & $26.6(6)^\circ$ & $f_2=0.168[0.160,0.175]$
      & $+0.0518[+0.00244,+0.1067]$ & $0.988[0.092,3.555]$\\
    $B_s$ & $\Bsone(5830)\to\Bs\gamma$ & legacy $\chi_s$
      & $38.5(1)^\circ$ & $f_2^{(B_s)}=0.0894[0.0854,0.0929]$
      & $-0.269[-0.302,-0.235]$ & $15.37[11.75,19.31]$\\
    $B_s$ & $\Bsone(5830)\to\Bs\gamma$ & lattice $f_{\gamma,s}^{\perp}$
      & $38.5(1)^\circ$ & $f_2^{(B_s)}=0.0888[0.0845,0.0926]$
      & $-0.203[-0.230,-0.175]$ & $8.702[6.474,11.23]$\\
    $B_s$ & $\Bsone^{\rm low}(5750)\to\Bs\gamma$ & legacy $\chi_s$
      & $38.5(1)^\circ$ & $f_1^{(B_s)}=0.533[0.508,0.557]$
      & $+0.216[+0.181,+0.249]$ & $5.435[3.949,7.588]$\\
    $B_s$ & $\Bsone^{\rm low}(5750)\to\Bs\gamma$ & lattice $f_{\gamma,s}^{\perp}$
      & $38.5(1)^\circ$ & $f_1^{(B_s)}=0.536[0.510,0.558]$
      & $+0.308[+0.288,+0.332]$ & $11.34[9.144,14.06]$\\
    \bottomrule
  \end{tabular}%
  }
  \caption{Final form-factor and width predictions.  The lattice photon input
  is the preferred modern normalization; the legacy susceptibility input is
  kept as a comparison with older photon-DA analyses.  The
  $\Bsone^{\rm low}(5750)$ entries are lower-partner diagnostics, not observed
  resonances.  Both sectors use direct $AA/AB/BB$ two-point residues.}
  \label{tab:results-polished}
\end{table}

The preferred lattice-normalized result for $\Dsone(2460)$ is a tens-of-keV
partial width.  In contrast, $\Dsone(2536)$ is highly cancellation sensitive:
the same axial and tensor basis amplitudes enter with the opposite physical
combination in Eq.~\eqref{eq:stage3}, driving the amplitude close to zero for a
wide part of the accepted input range.  The observed $\Bsone(5830)$ behaves
similarly in the bottom sector.  Its median width is at the keV scale after the
direct physical-current normalization.

Table~\ref{tab:interference-polished} makes this cancellation explicit.  The
entries are the two normalized pieces of Eq.~\eqref{eq:stage3}.  For
$\Dsone(2460)$ both pieces have the same sign and add constructively.  For
$\Dsone(2536)$ they have opposite signs and nearly cancel.  This is the LCSR
version of the Bondar--Milstein explanation of why the $\Dsone(2536)$ radiative
mode is difficult to observe.

\begin{table}[htbp]
  \centering
  \resizebox{\textwidth}{!}{%
  \begin{tabular}{llllccccc}
    \toprule
    Case & Mixing & $A_{2460}$ & $B_{2460}$ & $G_{2460}$
    & $A_{2536}$ & $B_{2536}$ & $G_{2536}$ & $|G_{2460}/G_{2536}|$\\
    \midrule
    previous-study nominal & $26.6^\circ$
      & $-0.0749$ & $-0.3365$ & $-0.4113$
      & $-0.3648$ & $+0.4110$ & $+0.0463$ & $8.89$\\
    truncated-matrix diagnostic & $30.61^\circ$
      & $-0.0850$ & $-0.3234$ & $-0.4084$
      & $-0.3597$ & $+0.4789$ & $+0.1192$ & $3.43$\\
    HQET benchmark & $35.3^\circ$
      & $-0.0966$ & $-0.3070$ & $-0.4035$
      & $-0.3358$ & $+0.5350$ & $+0.1992$ & $2.03$\\
    \bottomrule
  \end{tabular}%
  }
  \caption{Central lattice-photon benchmark at the midpoint of the selected
  transition window.  The $2460$ pieces remain constructive and the $2536$
  pieces remain destructive; the angle changes the distance from the
  cancellation near $23.96^\circ$.  All $G$ values are in $\GeV^{-1}$.}
  \label{tab:interference-polished}
\end{table}

The suppression is not a phase-space effect or an exact selection rule.  At
the nominal angle the two high-state pieces are individually of order
$0.4~\GeV^{-1}$,
\begin{equation}
G_{2536}^{(A)}=-0.3648~\GeV^{-1},\qquad
G_{2536}^{(B)}=+0.4110~\GeV^{-1},
\end{equation}
but their sum is only $+0.0463~\GeV^{-1}$.  Since the width is proportional
to $|G|^2q_\gamma^3$, this amplitude-level compensation becomes a much stronger
width suppression.  It is also why calculations with a different mixing
angle, current normalization, or relative axial--tensor basis amplitude can
obtain substantially larger widths: at the HQET benchmark angle the same
central OPE gives $G_{2536}=0.199~\GeV^{-1}$ and
$\Gamma_{2536}=12.3~\keV$.  The robust prediction is therefore the
destructive-interference mechanism and the hierarchy relative to the lower
state; the precise sub-keV median is intrinsically cancellation sensitive.

The lower $\Bsone$ row in Table~\ref{tab:results-polished} is retained only as
a model-state diagnostic.  It is not an experimentally established resonance
and is not included in the final experiment--theory comparison.

\begin{figure}[htbp]
  \centering
  \includegraphics[width=0.82\textwidth]{../outputs/final_window_mc_width_histograms.pdf}
  \caption{Final-window Monte Carlo width distributions with the preferred
  lattice photon normalization and the sector-specific previous-study mixing
  angles.}
  \label{fig:mc-polished}
\end{figure}

\begin{figure}[htbp]
  \centering
  \includegraphics[width=0.96\textwidth]{../outputs/final_angle_stability.pdf}
  \caption{Transition-LCSR stability at the previous-study angles
  $\theta_{D_s}=26.6\pm0.6^\circ$ and
  $\theta_{B_s}=38.5\pm0.1^\circ$.  The curves use the central angles and the
  shaded bands show only their quoted uncertainties.  The physical residues
  are held at their accepted-sample medians to isolate Borel and threshold
  dependence.  The high-state $D_s$ channel is visibly more sensitive because
  its axial and tensor contributions nearly cancel.}
  \label{fig:stability-polished}
\end{figure}

\section{Reproducibility}

Table~\ref{tab:repro-polished} records the source of each headline table and
figure.  The machine-readable version is
\texttt{outputs/manuscript\_reproducibility\_map.csv}.

\begin{table}[htbp]
  \centering
  \resizebox{\textwidth}{!}{%
  \begin{tabular}{llll}
    \toprule
    Manuscript item & Source script & Source output & Purpose\\
    \midrule
    input table & manual bookkeeping & \texttt{inputs\_table.csv} & input values and citations\\
    residue benchmark & \texttt{stage3\_residue\_benchmark\_checks.py} & \texttt{stage3\_residue\_benchmark\_checks.csv} & pure-current validation\\
    main result table & Stage-3 charm and bottom scripts & \texttt{combined\_recommended\_results\_table.csv} & final $G$ and $\Gamma$ values\\
    final comparison table & paper summary scripts & \texttt{bondar\_literature\_comparison\_table.csv}, \texttt{experimental\_comparison\_table.csv} & experiment and literature context\\
    MC width plots & lattice/physical-normalization scan scripts & \texttt{lattice\_fperp\_...pdf}, \texttt{stage3\_bs1\_physical\_...pdf} & uncertainty distributions\\
    final-angle stability & \texttt{final\_angle\_stability\_plots.py} & \texttt{final\_angle\_stability.csv}, \texttt{final\_angle\_stability.pdf} & Borel, threshold, and quoted-angle checks\\
    validation checks & validation and Ward scripts & \texttt{validation\_checks\_summary.txt}, \texttt{ward\_transversality\_check.txt} & limiting cases and gauge structure\\
    interference diagnostic & \texttt{ds1\_interference\_diagnostic.py} & \texttt{ds1\_interference\_diagnostic\_summary.csv} & constructive/destructive amplitude pieces\\
    \bottomrule
  \end{tabular}%
  }
  \caption{Reproducibility map for the headline material in this draft.}
  \label{tab:repro-polished}
\end{table}

\section{Conclusions}

We have developed a staged light-cone QCD sum-rule framework for the radiative
transitions of strange heavy axial mesons.  The calculation includes hard photon
emission from both quark lines, two-particle and three-particle photon
distribution amplitudes, and explicit strange-quark-mass terms.  The charm
residue normalization has been separated from the basis-current OPE and checked
by projection onto the pure axial channel.  This projection is the meaningful
comparison with Wang's one-current calculation; reproducing Wang's one-pole
number is not used as a criterion for the correctness of the two mixed
residues.  The bottom residues are now obtained from the direct heavy-flavour
continuation of the same normalized-current $AA/AB/BB$ matrix, rather than
from basis-current decay constants or an overlap closure.
The former QCD-side pole-mass numerator shortcut has been withdrawn.  The
complete off-shell double-Borel reduction given in the appendix has now been
propagated through the numerical uncertainty scan.  At the central charm and
bottom checkpoints, 145 convolution, invariant, coupling, residue, and width
quantities agree between the Python and Mathematica implementations within the
stated $10^{-8}$ absolute-or-relative regression tolerance.

Table~\ref{tab:conclusion-comparison} gives the requested consolidated
comparison with experiment and representative calculations.  The
$\Dsone(2460)$ interval overlaps the upper edge of the original single-current
LCSR range of Colangelo et al.\ \cite{Colangelo2005}.  Combining our preferred partial width
with the PDG 2025-update branching fraction
$\mathcal B(\Dsone(2460)\to\Ds\gamma)=18\pm4\%$ gives an indicative total width
near $0.19~\MeV$, comfortably below the direct $3.5~\MeV$ upper limit
\cite{PDG2025Update}.  For $\Dsone(2536)$, the measured
$0.92\pm0.05~\MeV$ total width implies
$\mathcal B_\gamma\simeq0.107\%$ at our median, with a 16--84 percentile range
of approximately $0.010$--$0.386\%$.  This is below the indirect
$\mathcal B_\gamma<0.9\%$, or $\Gamma_\gamma<8~\keV$, constraint inferred by
Bondar \cite{Bondar2025Why}.  The observed $\Bsone(5830)$ has no direct
radiative measurement, so that row remains a genuine prediction of the
direct-normalization calculation.

\begin{table}[htbp]
  \centering
  \footnotesize
  \setlength{\tabcolsep}{3pt}
  \begin{tabular}{L{0.24\textwidth}L{0.24\textwidth}
                  L{0.29\textwidth}L{0.16\textwidth}}
    \toprule
    Channel and this work [keV] & Experimental information
      & Selected theory [keV] & Assessment\\
    \midrule
    $\Dsone(2460)\to\Ds\gamma$\newline
      $34.96[28.06,43.53]$
      & $\mathcal B_\gamma=18(4)\%$;\newline
        $\Gamma_{\rm tot}<3.5~\MeV$
        \cite{PDG2025Update}
      & LCSR $19$--$29$ \cite{Colangelo2005};
        QM $3.53$--$17.2$ \cite{GoityRoberts2001,Chen2020Electric}
      & overlaps the upper edge of the earlier LCSR\\
    $\Dsone(2536)\to\Ds\gamma$\newline
      $0.988[0.092,3.555]$
      & $\Gamma_{\rm tot}=0.92(5)~\MeV$ \cite{PDG2025Update};\newline
        indirect $\Gamma_\gamma<8$ \cite{Bondar2025Why}
      & $1.6\pm2.3$ \cite{KornerPirjolSchilcher1993};
        $7$--$61$ (models, incl. $15$ from Ref.~\cite{Godfrey2003Radiative});
        $12$ \cite{BondarMilstein2025}
      & compatible with the bound; cancellation sensitive\\
    $\Bsone(5830)\to\Bs\gamma$\newline
      $8.702[6.474,11.23]$
      & no direct radiative measurement \cite{CMS2018}
      & $27$--$56.6$ \cite{LiNiZhong2021,GodfreyMoatsSwanson2016,LuPanWang2016};
        $28$ \cite{BondarMilstein2025}
      & below existing models; cancellation sensitive\\
    \bottomrule
  \end{tabular}
  \caption{Consolidated experiment and literature comparison for the preferred
  lattice-photon scenario.  The wide theory spread reflects different
  state assignments, mixing conventions, and normalization schemes.}
  \label{tab:conclusion-comparison}
\end{table}

The scientifically robust conclusion is the interference pattern rather than
an artificially precise high-state number.  The $\Dsone(2460)$ amplitude is
constructive and naturally lies at the tens-of-keV level.  The
$\Dsone(2536)$ amplitude is destructive and the final-angle stability plot
shows enhanced threshold and angle sensitivity near its zero.  The observed
$\Bsone(5830)$ is likewise a cancellation channel.  Its direct two-point
normalization raises the preferred width relative to the former closure
diagnostic, but it remains below the representative quark-model predictions.

\appendix

\section{Explicit analytic light-cone sum rules}
\label{app:lcsr}

\subsection{Correlator, Wick contractions, and propagators}

With $p'=p+q$, the Rohrwild-style external-photon correlator used for either
basis current is
\begin{equation}
 \Pi_\mu^X(p,q)=\ii\int d^4x\,e^{\ii p\cdot x}
 \langle\gamma(q,\varepsilon)|
 T\{J_\mu^X(x)J_5^\dagger(0)\}|0\rangle,\qquad X=A,B .
 \label{eq:app-correlator}
\end{equation}
The hard-emission Wick contraction is
\begin{equation}
 \Pi_{\mu}^{X,\mathrm{hard}}
 =-\ii N_c\!\int\!\frac{d^4k}{(2\pi)^4}
 \operatorname{Tr}\!\left[
 \Gamma_\mu^X S_Q^\gamma(k)\,\ii\gamma_5 S_s^0(k-p)
 +\Gamma_\mu^X S_Q^0(k)\,\ii\gamma_5 S_s^\gamma(k-p)
 \right],
 \label{eq:app-wick}
\end{equation}
where $S_f^0(k)=\ii(\slashed{k}+m_f)/(k^2-m_f^2+\ii0)$.
The light-strange propagator is written without ordinary local-condensate
pieces:
\begin{align}
 \widehat{s^a(x)\bar s^b(0)}={}&
 \delta^{ab}\left[\frac{\ii\slashed{x}}{2\pi^2x^4}
 -\frac{m_s}{4\pi^2x^2}\right]\nonumber\\
 &-\frac{\ii g_s}{16\pi^2x^2}\int_0^1du\,
 \{\bar u\slashed{x}\sigma_{\alpha\beta}
 +u\sigma_{\alpha\beta}\slashed{x}\}
 G_A^{\alpha\beta}(ux)t^A_{ab}\nonumber\\
 &-\frac{\ii e_s\delta^{ab}}{16\pi^2x^2}\int_0^1du\,
 \{\bar u\slashed{x}\sigma_{\alpha\beta}
 +u\sigma_{\alpha\beta}\slashed{x}\}
 F^{\alpha\beta}(ux)\nonumber\\
 &-\frac{m_s}{32\pi^2}\int_0^1du\,
 [g_sG_A^{\alpha\beta}(ux)t^A_{ab}
 +e_s\delta^{ab}F^{\alpha\beta}(ux)]
 \sigma_{\alpha\beta}L_x+\cdots ,
 \label{eq:app-light-propagator}
\end{align}
where $\bar u=1-u$ and
$L_x=\ln(-x^2\mu^2/4)+2\gamma_E$.  The last line is the linear-$m_s$ Wilson
coefficient in one background field; it is not a local condensate.  Setting
$m_s=0$ reproduces the form used by Rohrwild \cite{Rohrwild2007}.

The one-background-field momentum representation used for both heavy and
light lines is
\begin{align}
 S_f^{\mathcal F}(x)=-\ii c_f\int\frac{d^4k}{(2\pi)^4}e^{-\ii k\cdot x}
 \int_0^1du\Bigg[
 &\frac{\slashed{k}+m_f}{2(m_f^2-k^2)^2}
 \mathcal F^{\alpha\beta}(ux)\sigma_{\alpha\beta}\nonumber\\
 &+\frac{ux_\alpha}{m_f^2-k^2}
 \mathcal F^{\alpha\beta}(ux)\gamma_\beta\Bigg],
 \label{eq:app-background-propagator}
\end{align}
with $(c_f,\mathcal F)=(e_f,F)$ for the electromagnetic field and
$(g_s,G)$ for the gluon field.  Equations
\eqref{eq:app-light-propagator}--\eqref{eq:app-background-propagator} show
explicitly that no ordinary local term such as
$-\langle\bar ss\rangle/12$ is inserted into the transition calculation.
The condensate-dependent terms instead occur in nonlocal photon matrix
elements.  In Rohrwild notation,
\begin{align}
 \langle0|\bar s(0)[0,x]\sigma_{\mu\nu}s(x)|0\rangle_F
 ={}&e e_s\ssbar\int_0^1du\Bigg\{
 \left[\chi_\gamma\phi_\gamma(u)+\frac{x^2}{16}\mathcal A(u)\right]
 F_{\mu\nu}(ux)\nonumber\\
 &+\frac{\mathcal B(u)}8x^\rho
 [x_\nu F_{\mu\rho}(ux)-x_\mu F_{\nu\rho}(ux)]
 \Bigg\}.
 \label{eq:app-photon-matrix-element}
\end{align}
The operator dictionary to the alternative notation is
\begin{equation}
 \mathcal B(u)=-4\int_0^u d\alpha\,(u-\alpha)h_\gamma(\alpha),
 \qquad
 \mathbb A(u)-8\bar H_\gamma(u)=\mathcal A(u)+2\mathcal B(u).
 \label{eq:app-Rohrwild-B-dictionary}
\end{equation}
In particular, the Rohrwild counterpart of the Colangelo
$\bar H_\gamma$ term is $2\mathcal B$, not a new $H_\gamma$ amplitude.

\subsection{Double-Borel normalization and axial invariant}

Define
\begin{align}
 \langle0|J_5|P(p)\rangle&=\frac{f_Pm_P^2}{m_Q+m_s},&
 \langle0|J_i^\mu|H_i(p',\eta)\rangle&=f_iM_i\eta^\mu .
\end{align}
After projection onto
$T_{\mu\nu}=(p\!\cdot q)g_{\mu\nu}-q_\mu p_\nu$, the pole contribution is
\begin{equation}
 \widehat T_i^{\rm pole}
 =\frac{f_Pm_P^2}{m_Q+m_s}\,f_iM_iG_i
 \exp\!\left(-\frac{M_i^2}{\mathcal M_1^2}
             -\frac{m_P^2}{\mathcal M_2^2}\right).
 \label{eq:app-pole}
\end{equation}
For $\mathcal M_1^2=\mathcal M_2^2=2M^2$,
\begin{align}
 G_1&=\mathcal N_1(\sin\theta\,\widehat T_A
                   +\cos\theta\,\widehat T_B),\nonumber\\
 G_2&=\mathcal N_2(\cos\theta\,\widehat T_A
                   -\sin\theta\,\widehat T_B),\qquad
 \mathcal N_i=
 \frac{e^{(M_i^2+m_P^2)/(2M^2)}(m_Q+m_s)}
 {M_if_i\,m_P^2f_P}.
 \label{eq:app-physical-lcsr}
\end{align}

Let $E_Q=e^{-m_Q^2/M^2}$, $E_0=e^{-s_0/M^2}$,
$\Delta E_Q=E_Q-E_0$, and $u_0=1/2$.  In the
positive-magnitude susceptibility convention used in the numerical code, the
paper-default axial invariant is
\begin{align}
 \widehat T_A^{\rm RW}={}&
 \int_{(m_Q+m_s)^2}^{s_0}ds\,e^{-s/M^2}\rho_A^P(s)
 +e_s\ssbar\Delta E_QM^2\chi_\gamma\phi_\gamma(u_0)\nonumber\\
 &+e_s\ssbar E_Q\left[
 \frac14\{\mathcal A(u_0)+2\mathcal B(u_0)\}
 \left(1+\frac{m_Q^2}{M^2}\right)
 +\mathcal I_{\mathcal F_1}(u_0)\right]\nonumber\\
 &+e_sf_{3\gamma}m_QE_Q\overline\psi^{(V)}(u_0)
 +\widehat T_A^{\rm 3p,\gamma}.
 \label{eq:app-axial-invariant}
\end{align}
Thus the continuum factors are not universal:
twist~2 carries $\Delta E_Q$, whereas twist~3, twist~4, and the gluonic
three-particle term carry $E_Q$.  The separate gauge-completion
electromagnetic term below follows Rohrwild's $I_F$-type
$\Delta E_Q$ subtraction.
Here
$\overline\psi^{(V)}(u)=2\int_0^u dv\,\psi^v(v)$, and
\begin{align}
 \mathcal I_{\mathcal F_1}(u_0)={}&
 \int_0^{1-u_0}\!dv\int_0^{u_0/(1-v)}\!d\alpha_g\,
 \mathcal F_1(u_0-(1-v)\alpha_g,1-u_0-v\alpha_g,\alpha_g)
 \nonumber\\
 &+\int_{1-u_0}^{1}\!dv\int_0^{(1-u_0)/v}\!d\alpha_g\,
 \mathcal F_1(u_0-(1-v)\alpha_g,1-u_0-v\alpha_g,\alpha_g).
 \label{eq:app-f1-integral}
\end{align}
The three-particle combination appearing here is
\begin{equation}
 \mathcal F_1=\mathcal S+\widetilde{\mathcal S}
 -\mathcal T_1-\mathcal T_2+\mathcal T_3+\mathcal T_4
 +2v(-\mathcal S-\mathcal T_3+\mathcal T_2).
\end{equation}
The perturbative density is unambiguously specified by
\begin{align}
 \rho_i(s;m_i,m_j,e_i)={}&-\frac{3e_i}{8\pi^2}
 \Bigg[2m_i\ln\frac{s-m_j^2+m_i^2-\sqrt{\lambda(s,m_j^2,m_i^2)}}
 {s-m_j^2+m_i^2+\sqrt{\lambda(s,m_j^2,m_i^2)}}\nonumber\\
 &+(m_j-m_i)\frac{m_j^2-m_i^2-s}{s^2}
 \sqrt{\lambda(s,m_j^2,m_i^2)}\Bigg],\nonumber\\
 \rho_A^P(s)={}&\rho_s(s;m_s,m_Q,e_s)-\rho_Q(s;m_Q,m_s,e_Q).
 \label{eq:app-rho-a}
\end{align}

\subsection{Tensor invariant and local-scheme separation}

For completeness, the tensor hard density used in the diagonal
double-dispersion prescription can be written
\begin{align}
 \rho_B^P(s)&=e_s\mathcal L(s;m_s,m_Q;a_s,b_s)
              -e_Q\mathcal L(s;m_Q,m_s;a_Q,b_Q),\nonumber\\
 \mathcal L(s;m_i,m_j;a,b)&=-\frac{3}{8\pi^2}\left[
 2a\ln\frac{s-m_j^2+m_i^2-\sqrt{\lambda}}
 {s-m_j^2+m_i^2+\sqrt{\lambda}}
 +\frac{b}{m_i^2}\frac{m_j^2-m_i^2-s}{s^2}\sqrt{\lambda}\right],
 \label{eq:app-rho-b}
\end{align}
where $d=m_Q+m_s$,
\begin{equation}
 a_s=\frac{m_Qm_s}{d},\quad b_s=\frac{m_s^2s}{d},\qquad
 a_Q=\frac{2m_Q^2-m_Qm_s}{d},\quad b_Q=-\frac{m_Q^2s}{d}.
\end{equation}
On the QCD side, $p'^2$ and $p^2$ remain independent until the double Borel
transformation:
\begin{equation}
 q^2=0,\qquad p'=p+q,\qquad
 2p\cdot q=p'^2-p^2,\qquad
 D(u)=m_Q^2-\bar u p'^2-up^2 .
 \label{eq:app-offshell-kinematics}
\end{equation}
The substitutions $p'^2=M_i^2$, $p^2=m_P^2$ and
$2p\cdot q=M_i^2-m_P^2$ belong only to the isolated phenomenological pole
residue and must not be made inside a QCD-side numerator.

The current traces give the exact kinematics-independent relations
\begin{equation}
 \widehat T_B^{\rm tw2}
 =\frac{m_Q}{d}\widehat T_A^{\rm tw2},\qquad
 \widehat T_B^{\rm tw4}
 =\frac{m_Q}{d}\widehat T_A^{\rm tw4}.
 \label{eq:app-tensor-safe-ratios}
\end{equation}
The vector-DA and $\mathbb P$ three-particle terms contain
$p^2$ or $p\cdot q$ and must therefore be transformed off shell.  For
$I_n[F]=\int_0^1du\,F(u)/D(u)^n$, define
\begin{equation}
 \frac1{M^2}=\frac1{\mathcal M_1^2}+\frac1{\mathcal M_2^2},\qquad
 u_0=\frac{\mathcal M_1^2}{\mathcal M_1^2+\mathcal M_2^2},\qquad
 A_n=\frac{(M^2)^{2-n}}{\Gamma(n)}e^{-m_Q^2/M^2}.
\end{equation}
Then, for $n>1$ and vanishing endpoint terms,
\begin{align}
 \mathcal B_{12} I_n[F]&=A_nF(u_0),\nonumber\\
 \mathcal B_{12}\!\int_0^1du\,\frac{(p\cdot q)F(u)}{D(u)^n}
 &=\frac{A_{n-1}}{2(n-1)}F'(u_0),\nonumber\\
 \mathcal B_{12}\!\int_0^1du\,\frac{p^2F(u)}{D(u)^n}
 &=m_Q^2A_nF(u_0)-A_{n-1}F(u_0)
 -\frac{A_{n-1}}{n-1}[\bar uF(u)]'_{u=u_0}.
 \label{eq:app-borel-numerators}
\end{align}
Nonvanishing endpoints generate additional contact terms.  These identities
replace the former pole-mass numerator prescription.

For the vector DA set $a=1-u$, $a_0=1-u_0$ and
$D_a=m_Q^2-ap'^2-(1-a)p^2$.  The reduced pre-Borel kernels and the explicit
tensor-current Borel result are
\begin{align}
 K_A^{(V)}(a)&=-\frac{8m_Q}{D_a^2},\nonumber\\
 K_B^{(V)}(a)&=-\frac{8[m_Q^2+(1-a)p\cdot q]}{dD_a^2},
 \nonumber\\[-2pt]
 \widehat T_B^{\rm tw3}
 &=\frac{e_sf_{3\gamma}E_Q}{d}
 \left[
 m_Q^2\overline\psi^{(V)}(a_0)
 -\frac{M^2}{2}
 \left.\frac{d}{da}
 \{(1-a)\overline\psi^{(V)}(a)\}\right|_{a=a_0}
 \right].
 \label{eq:app-tensor-tw3-explicit}
\end{align}
At $a_0=1/2$,
\begin{equation}
 \widehat T_A^{\rm tw3}=0,\qquad
 \widehat T_B^{\rm tw3}
 =-\frac{e_sf_{3\gamma}M^2E_Q}{4d}
 \overline\psi^{(V)\prime}(1/2).
 \label{eq:app-symmetric-tw3}
\end{equation}

For the three-particle pieces the heavy-line fraction is
$a(\alpha_i,v)=\alpha_{\bar q}+v\alpha_g$.  After performing the
delta-function integration over the physical simplex, define
\begin{align}
 \mathcal J_z[\mathcal F]={}&
 \int_0^z dv\int_0^{(1-z)/(1-v)}d\alpha_g\,
 \mathcal F(1-z-(1-v)\alpha_g,z-v\alpha_g,\alpha_g,v)
 \nonumber\\
 &+\int_z^1dv\int_0^{z/v}d\alpha_g\,
 \mathcal F(1-z-(1-v)\alpha_g,z-v\alpha_g,\alpha_g,v).
 \label{eq:app-line-projection}
\end{align}
The $1/(q\cdot x)$ terms require the support density
\begin{align}
 \mathcal L_z[X]={}&
 \int_0^{1-z}d\alpha_q\int_0^{\alpha_q}d\alpha_q'
 \int_0^z d\alpha_{\bar q}\,
 \frac{z-\alpha_{\bar q}}{(1-\alpha_q-\alpha_{\bar q})^2}
 X(\alpha_q,\alpha_{\bar q},1-\alpha_q-\alpha_{\bar q})
 \nonumber\\
 &-\int_0^{1-z}d\alpha_q\int_z^{1-\alpha_q}d\beta\,
 \frac{z}{\beta^2}\int_0^\beta d\alpha_{\bar q}'\,
 X(\alpha_q,\alpha_{\bar q}',1-\alpha_q-\alpha_{\bar q}')
 \nonumber\\
 &+\int_0^{1-z}d\alpha_q\,\frac{z}{(1-\alpha_q)^2}
 \int_0^{\alpha_q}d\alpha_q'
 \int_0^{1-\alpha_q'}d\alpha_{\bar q}\,
 X(\alpha_q',\alpha_{\bar q},1-\alpha_q'-\alpha_{\bar q}).
 \label{eq:app-P-line-density}
\end{align}
The lower limit $\beta=z$ in the second line is required by the double-Borel
support; extending it to zero produces a spurious logarithmic divergence.

The $x_\alpha G^{\alpha\beta}\gamma_\beta$ heavy-propagator term has no
$(\slashed k+m_Q)$ numerator.  Its explicit E1 trace therefore gives
\begin{align}
 \mathfrak X_{A,S}[F]&=-8E_QF(a_0),&
 \mathfrak X_{A,T_1}[F]&=48\ii E_QF(a_0),\nonumber\\
 \mathfrak X_{A,T_2}[F]&=16\ii E_QF(a_0),&
 \mathfrak X_{A,T_3}[F]&=-16\ii E_QF(a_0),\nonumber\\
 \mathfrak X_{A,T_4}[F]&=0,
 \label{eq:app-A-xgamma-explicit}\\
 \mathfrak X_{B,S}[F]&=\mathfrak X_{B,T_1}[F]
 =\mathfrak X_{B,T_2}[F]=\mathfrak X_{B,T_3}[F]
 =\mathfrak X_{B,T_4}[F]=0 .
 \label{eq:app-B-xgamma-explicit}
\end{align}
The corresponding $\sigma_{\alpha\beta}$ coefficients are
$\{8,-16\ii\}$ for $A$ and
$\{8m_Q/d,-16\ii m_Q/d\}$ for $B$.
Equations~\eqref{eq:app-A-xgamma-explicit}--\eqref{eq:app-B-xgamma-explicit}
are raw Borel images.  The final continuum prescription is term-specific:
twist~2 uses $\Delta E_Q$, twist~3, twist~4, and the gluonic
three-particle term retain $E_Q$, and the separate gauge-completion
electromagnetic term uses $\Delta E_Q$.

Suppressing the arguments $(\alpha_q,\alpha_{\bar q},\alpha_g)$, the complete
gluonic operator weights are
\begin{align}
 \mathcal F^g_{A,\sigma}
 &=\mathcal S+\widetilde{\mathcal S}
   -\mathcal T_1-\mathcal T_2+\mathcal T_3+\mathcal T_4,\nonumber\\
 \mathcal F^g_{A,x\gamma}
 &=2v(-\mathcal S-\mathcal T_3+\mathcal T_2),\nonumber\\
 \mathcal F_A^g&=\mathcal F^g_{A,\sigma}
                 +\mathcal F^g_{A,x\gamma},&
 \mathcal F_B^g&=\frac{m_Q}{d}\mathcal F^g_{A,\sigma}.
 \label{eq:app-gluonic-weights}
\end{align}
Their explicit post-Borel invariants are
\begin{equation}
 \widehat T_A^{\rm 3p,g}
 =e_s\ssbar E_Q\mathcal J_{a_0}[\mathcal F_A^g],\qquad
 \widehat T_B^{\rm 3p,g}
 =e_s\ssbar E_Q\mathcal J_{a_0}[\mathcal F_B^g].
 \label{eq:app-gluonic-terms}
\end{equation}

For the nonlocal electromagnetic operators,
\begin{equation}
 \mathcal F_A^\gamma=(1-2v)\mathcal S_\gamma-\mathcal T_4^\gamma,
 \qquad
 \mathcal F_B^\gamma=\frac{m_Q}{d}
 (\mathcal S_\gamma-\mathcal T_4^\gamma).
 \label{eq:app-em-weights}
\end{equation}
The exact off-shell Borel transform replaces the former pole-mass numerator
by
\begin{equation}
 \frac{2(m_P^2-M_i^2)}{M^2}\,
 \mathcal L_{a_0}[\mathcal T_4^\gamma]
 \longrightarrow
 2\left.\frac{d}{dz}\mathcal L_z[\mathcal T_4^\gamma]
 \right|_{z=a_0}.
 \label{eq:app-P-borel-replacement}
\end{equation}
Writing $\Delta E_Q=E_Q-E_0$, the complete electromagnetic terms are
\begin{align}
 \widehat T_A^{\rm 3p,\gamma}
 &=e_Q\ssbar\Delta E_Q
 \left\{\mathcal J_{a_0}[
 (1-2v)\mathcal S_\gamma-\mathcal T_4^\gamma]
 +2\left.\frac{d}{dz}\mathcal L_z[\mathcal T_4^\gamma]
 \right|_{z=a_0}\right\},\nonumber\\
 \widehat T_B^{\rm 3p,\gamma}
 &=e_Q\ssbar\Delta E_Q\frac{m_Q}{d}
 \left\{\mathcal J_{a_0}[
 \mathcal S_\gamma-\mathcal T_4^\gamma]
 +2\left.\frac{d}{dz}\mathcal L_z[\mathcal T_4^\gamma]
 \right|_{z=a_0}\right\}.
 \label{eq:app-em-terms}
\end{align}
Thus the Rohrwild tensor invariant has the parallel structure
\begin{equation}
 \widehat T_B^{\rm RW}=\int ds\,e^{-s/M^2}\rho_B^P
 +\widehat T_B^{\rm tw2}
 +\widehat T_B^{\rm tw3}
 +\widehat T_B^{\rm tw4}
 +\widehat T_B^{\rm 3p,g}
 +\widehat T_B^{\rm 3p,\gamma}.
\end{equation}
Equations~\eqref{eq:app-line-projection}--\eqref{eq:app-em-terms} display all
integration limits, operator weights and Borel derivatives.  No physical
pole mass and no ordinary local condensate enters a QCD-side numerator.

\subsection{Numerical convolution and cross-code checkpoint}

At the symmetric saddle the DA-only convolutions are
\begin{align}
 \mathcal J_{1/2}[\mathcal F^g_{A,\sigma}]&=-0.1562500000,&
 \mathcal J_{1/2}[\mathcal F^g_{A,x\gamma}]&=+0.2265798500,\nonumber\\
 \mathcal J_{1/2}[\mathcal F^g_A]&=+0.0703298500,&
 \mathcal J_{1/2}[\mathcal F^\gamma_A]&=0,\nonumber\\
 \mathcal J_{1/2}[\mathcal S_\gamma-\mathcal T_4^\gamma]
 &=1.5625000000,&
 \mathcal L_{1/2}[\mathcal T_4^\gamma]&=0.1111834583,\nonumber\\
 \left.\frac{d}{dz}\mathcal L_z[\mathcal T_4^\gamma]\right|_{z=1/2}
 &=0.4098669166.&&
 \label{eq:app-numerical-convolutions}
\end{align}
For the central charm regression point
$M^2=3.75~\GeV^2$, $s_0=8.0~\GeV^2$, the lattice photon
normalization, $\theta=26.6^\circ$, and
$(f_1,f_2)=(0.4005893,0.1666704)~\GeV$, the invariant breakdown is
\begin{table}[H]
 \centering
 \begin{tabular}{lrr}
  \toprule
  contribution & $\widehat T_A~(\GeV^3)$ & $\widehat T_B~(\GeV^3)$\\
  \midrule
  perturbative & $-8.8661701\times10^{-2}$ & $-1.1226165\times10^{-1}$\\
  twist 2 & $+5.4942081\times10^{-2}$ & $+5.1193282\times10^{-2}$\\
  twist 3 & $0$ & $-5.2253296\times10^{-3}$\\
  twist 4 & $+5.5991083\times10^{-3}$ & $+5.2170708\times10^{-3}$\\
  three-particle gluonic & $+1.6863557\times10^{-4}$ & $-3.4908999\times10^{-4}$\\
  three-particle electromagnetic & $-3.2152535\times10^{-3}$ & $-8.7063197\times10^{-3}$\\
  \midrule
  total & $-3.1167129\times10^{-2}$ & $-7.0132040\times10^{-2}$\\
  \bottomrule
 \end{tabular}
 \caption{Central invariant checkpoint.  The corresponding single-point
 couplings are $G_1=-0.4113424~\GeV^{-1}$ and
 $G_2=+0.0464987~\GeV^{-1}$, giving $35.5665$ and $0.67094~\keV$.
 These are regression values, not the Monte Carlo medians in
 Table~\ref{tab:results-polished}.}
 \label{tab:app-central-invariant-check}
\end{table}
The independent Python and Mathematica tables contain 145 common numerical
entries and pass an absolute-or-relative $10^{-8}$ comparison.  The executable
sources and machine-readable residuals are, respectively,
\begin{center}
\small
\path{scripts/corrected_transition_central_python.py},\\
\path{scripts/mathematica_corrected_transition_numerics.wl},\\
\path{outputs/corrected_transition_python_mathematica_comparison.csv}.
\end{center}

\begin{table}[htbp]
 \centering
 \resizebox{\textwidth}{!}{%
 \begin{tabular}{lll}
  \toprule
  Contribution & Rohrwild default & Colangelo comparison\\
  \midrule
  ordinary local transition condensate & absent & retained for $J_A$\\
  $\mathcal S_\gamma,\mathcal T_4^\gamma$ & included nonlocally & not added separately\\
  gluonic photon DAs & included & included through $\mathcal F_1$\\
  local condensates in two-point $f_i$ sum rules & required & required\\
  local plus electromagnetic term together & forbidden double counting & forbidden double counting\\
  \bottomrule
 \end{tabular}
 }
 \caption{Separation of the transition and two-point condensate prescriptions.}
 \label{tab:app-local-schemes}
\end{table}

The alternative Colangelo axial invariant is obtained by removing
$\widehat T_A^{\rm 3p,\gamma}$ from Eq.~\eqref{eq:app-axial-invariant} and adding
\begin{equation}
 e_QE_Q\ssbar\left[1-\frac{m_Qm_s}{M^2}
 +\frac{m_s^2}{2M^2}\left(1+\frac{m_Q^2}{M^2}\right)\right].
 \label{eq:app-colangelo-local}
\end{equation}
The two terms are alternatives and are never summed.  In our rotation
convention the exact pure-axial Colangelo limit is
\begin{equation}
 \theta=\frac{\pi}{2},\qquad J_1=J_A,\qquad
 G_1=\frac{e^{(M_1^2+m_P^2)/(2M^2)}(m_Q+m_s)}
 {M_1f_A\,m_P^2f_P}\widehat T_A .
 \label{eq:app-colangelo-limit}
\end{equation}
Thus the agreement in analytic structure is exact at the pure-current point;
with the Colangelo inputs our numerical pure-$A$ width interval is
$20.51$--$33.26~\keV$, overlapping their $19$--$29~\keV$ result.

\section{Explicit \texorpdfstring{$AA/AB/BB$}{AA/AB/BB} two-point sum rules}
\label{app:two-point}

For $i,j\in\{A,B\}$,
\begin{equation}
 \Pi_{ij}^{\mu\nu}(p')=\ii\int d^4x\,e^{\ii p'\cdot x}
 \langle0|T\{J_i^\mu(x)J_j^{\nu\dagger}(0)\}|0\rangle
 =\left(-g^{\mu\nu}+\frac{p'^\mu p'^\nu}{p'^2}\right)
 \Pi_{ij}(p'^2)+\cdots .
\end{equation}
Let $d=m_Q+m_s$ and
$\lambda(s)=[s-(m_Q+m_s)^2][s-(m_Q-m_s)^2]$.  The exact-mass LO spectral
matrix is
\begin{align}
 \rho_{AA}(s)&=\frac{\sqrt\lambda}{8\pi^2}\left[
 2-\frac{m_Q^2+m_s^2+6m_Qm_s}{s}
 -\frac{(m_Q^2-m_s^2)^2}{s^2}\right],\nonumber\\
 \rho_{AB}(s)&=\rho_{BA}(s)=\frac{\sqrt\lambda}{8\pi^2}
 \frac{3(m_s-m_Q)(d^2-s)}{ds},\nonumber\\
 \rho_{BB}(s)&=\frac{\sqrt\lambda}{8\pi^2}
 \frac{-(d^2-s)(2m_s^2-4m_Qm_s+2m_Q^2+s)}{d^2s}.
 \label{eq:app-two-point-rho}
\end{align}
With $C_3=\ssbar$ and
$C_5=\langle\bar s g_s\sigma\!\cdot\!G s\rangle$, the common
LO+$d=3$+$d=5$ Borel matrix is
\begin{equation}
 \widehat\Pi(M^2,s_0)=
 \int_{d^2}^{s_0}ds\,e^{-s/M^2}\rho(s)
 +e^{-m_Q^2/M^2}
 [C_3(\mathcal Q_0+\mathcal Q_1+\mathcal Q_2)+C_5\mathcal Q_5],
 \label{eq:app-two-point-matrix}
\end{equation}
where
\begin{align}
 \mathcal Q_0&=
 \begin{pmatrix}m_Q&m_Q^2/d\\m_Q^2/d&m_Q^3/d^2\end{pmatrix},\nonumber\\
 \mathcal Q_1&=
 \begin{pmatrix}
 -\dfrac{m_sm_Q^2}{2M^2}&
 \dfrac{m_sm_Q(1-m_Q^2/M^2)}{2d}\\
 \dfrac{m_sm_Q(1-m_Q^2/M^2)}{2d}&
 \dfrac{m_sm_Q^2(1-m_Q^2/(2M^2))}{d^2}
 \end{pmatrix},\nonumber\\
 \mathcal Q_2&=
 \begin{pmatrix}
 \dfrac{m_s^2m_Q^3}{2M^4}&
 \dfrac{m_s^2(m_Q^4/M^4-m_Q^2/M^2-1)}{2d}\\
 \dfrac{m_s^2(m_Q^4/M^4-m_Q^2/M^2-1)}{2d}&
 \dfrac{m_s^2(m_Q^5/(2M^4)-m_Q^3/M^2)}{d^2}
 \end{pmatrix},\nonumber\\
 \mathcal Q_5&=
 \begin{pmatrix}
 -\dfrac{m_Q^3}{4M^4}&
 -\dfrac{m_Q^4/M^4-m_Q^2/M^2-1}{4d}\\
 -\dfrac{m_Q^4/M^4-m_Q^2/M^2-1}{4d}&
 -\dfrac{m_Q^5/(4M^4)-m_Q^3/(2M^2)}{d^2}
 \end{pmatrix}.
 \label{eq:app-q-matrices}
\end{align}
These local condensates are required because this is an ordinary two-point
SVZ sum rule; their use does not contradict their exclusion from the
external-photon transition LCSR.

Writing $s_\theta=\sin\theta$ and $c_\theta=\cos\theta$, the physical
invariants are
\begin{align}
 \widehat\Pi_{11}&=s_\theta^2\widehat\Pi_{AA}
 +c_\theta^2\widehat\Pi_{BB}
 +2s_\theta c_\theta\widehat\Pi_{AB},\nonumber\\
 \widehat\Pi_{22}&=c_\theta^2\widehat\Pi_{AA}
 +s_\theta^2\widehat\Pi_{BB}
 -2s_\theta c_\theta\widehat\Pi_{AB},\nonumber\\
 \widehat\Pi_{12}&=s_\theta c_\theta
 (\widehat\Pi_{AA}-\widehat\Pi_{BB})
 +(c_\theta^2-s_\theta^2)\widehat\Pi_{AB}.
 \label{eq:app-rotated-pi}
\end{align}
For $n=0,1,2,5$, define the fully specified scalar projections of the matrices
in Eq.~\eqref{eq:app-q-matrices} by
\begin{align}
 q_1^{(n)}&=s_\theta^2(\mathcal Q_n)_{AA}
 +c_\theta^2(\mathcal Q_n)_{BB}
 +2s_\theta c_\theta(\mathcal Q_n)_{AB},\nonumber\\
 q_2^{(n)}&=c_\theta^2(\mathcal Q_n)_{AA}
 +s_\theta^2(\mathcal Q_n)_{BB}
 -2s_\theta c_\theta(\mathcal Q_n)_{AB}.
\end{align}
The two pole channels use separate thresholds.  The expressions actually
evaluated are
\begin{align}
 \widehat\Pi_{11}^{(1)}={}&
 \int_{d^2}^{s_0^{(1)}}ds\,e^{-s/M^2}
 [s_\theta^2\rho_{AA}+c_\theta^2\rho_{BB}
  +2s_\theta c_\theta\rho_{AB}]
 \nonumber\\
 &+e^{-m_Q^2/M^2}
 \{C_3[q_1^{(0)}+q_1^{(1)}+q_1^{(2)}]+C_5q_1^{(5)}\},
 \nonumber\\
 \widehat\Pi_{22}^{(2)}={}&
 \int_{d^2}^{s_0^{(2)}}ds\,e^{-s/M^2}
 [c_\theta^2\rho_{AA}+s_\theta^2\rho_{BB}
  -2s_\theta c_\theta\rho_{AB}]
 \nonumber\\
 &+e^{-m_Q^2/M^2}
 \{C_3[q_2^{(0)}+q_2^{(1)}+q_2^{(2)}]+C_5q_2^{(5)}\}.
\end{align}
Consequently,
\begin{align}
 \boxed{f_1(M^2,s_0^{(1)},\theta)}
 &=\frac{e^{M_1^2/(2M^2)}}{M_1}
 \sqrt{\widehat\Pi_{11}^{(1)}},\nonumber\\
 \boxed{f_2(M^2,s_0^{(2)},\theta)}
 &=\frac{e^{M_2^2/(2M^2)}}{M_2}
 \sqrt{\widehat\Pi_{22}^{(2)}}.
 \label{eq:app-decay-constants}
\end{align}
We insert the previous-study angle $26.6\pm0.6^\circ$
\cite{AlievBilmisSavci2024} here; we do not refit it in the radiative paper.
Setting $\widehat\Pi_{12}=0$ gives the
$30.63[29.94,31.29]^\circ$ truncation diagnostic quoted in the text.
At the central point
$M^2=2.35~{\rm GeV}^2$,
$s_0^{(1)}=7.925~{\rm GeV}^2$ and
$s_0^{(2)}=9.975~{\rm GeV}^2$, these expressions give
$f_1=0.4005893~{\rm GeV}$ and $f_2=0.1666704~{\rm GeV}$.  The values quoted
in the results tables are the corresponding Monte Carlo medians and
16--84 percentile intervals.

For the bottom channel the same formulas are evaluated with $m_Q=m_b$,
$\theta_{B_s}=38.5^\circ$, $M^2=7.00~\GeV^2$,
$s_0^{(1)}=44.35~\GeV^2$, and
$s_0^{(2)}=43.025~\GeV^2$.  The central direct residues are
\begin{equation}
 f_1^{(B_s)}=0.5381107~\GeV,\qquad
 f_2^{(B_s)}=0.0888787~\GeV .
\end{equation}
The corresponding pole fractions are $0.805$ and $0.556$, while the
dimension-five fractions are $0.0029$ and $0.0348$.  Diagonalizing the
truncated bottom matrix gives $39.50^\circ$ as a diagnostic, close to the
external angle; the normalized off-diagonal residual is $0.104$.  Varying the
OPE inputs and the two-point Borel parameter produces the Monte Carlo
intervals quoted in the text.

\subsection{Projection onto Wang's pure axial channel}

Inverting the physical rotation gives $J_A=s_\theta J_1+c_\theta J_2$.
Consequently the two-pole hadronic representation of the pure-$AA$ invariant is
\begin{equation}
 \widehat\Pi_{AA}^{(2\,{\rm pole})}
 =s_\theta^2 f_1^2M_1^2e^{-M_1^2/M^2}
 +c_\theta^2 f_2^2M_2^2e^{-M_2^2/M^2}.
 \label{eq:app-aa-two-pole}
\end{equation}
For comparison with a one-pole result normalized at $M_1$, define only as a
diagnostic
\begin{equation}
 f_{A,\mathrm{eq}}^2(M^2)=
 \frac{e^{M_1^2/M^2}}{M_1^2}
 \widehat\Pi_{AA}^{(2\,{\rm pole})}.
 \label{eq:app-wang-projection}
\end{equation}
The accepted two-point samples give
$f_{A,\mathrm{eq}}=0.2303[0.2194,0.2418]~\GeV$, while the separate direct
$AA$ OPE in Wang's window gives
$0.2686[0.2607,0.2768]~\GeV$ and Wang quotes
$0.345\pm0.017~\GeV$ \cite{Wang2015}.  The three quantities test related
normalization scales but are not identical calculations: the first uses two
physical poles, the second uses our stated truncated OPE, and the third uses
Wang's one-pole OPE.  This is why the comparison is useful but cannot by itself
prove or disprove the mixed-current result.

\section*{Acknowledgments}

To be added.

\bibliographystyle{unsrt}
\bibliography{references}

\end{document}
